% Figure: RMS value of a sinusoid. Shows v(t) = Vp sin(wt), its square,
% and the horizontal line at the mean of the square, whose square root
% is the RMS value Vp/sqrt(2).
\begin{figure}[htbp]
\centering
\begin{tikzpicture}
\begin{axis}[
  width=10.5cm, height=6cm,
  xlabel={$t/T$}, ylabel={signal (normalised to $V_{p}$)},
  xmin=0, xmax=1, ymin=-1.15, ymax=1.15,
  domain=0:1, samples=200,
  axis lines=middle, grid=both,
  major grid style={engblack!12}, font=\small,
  legend pos=south east, legend style={font=\footnotesize},
]
\addplot[engblue, very thick] {sin(deg(2*pi*x))};
\addlegendentry{$v(t)/V_{p} = \sin(\omega t)$}
\addplot[engred, very thick] {sin(deg(2*pi*x))^2};
\addlegendentry{$v^{2}(t)/V_{p}^{2}$}
\addplot[engamber, very thick, dashed, domain=0:1] {0.5};
\addlegendentry{mean of $v^{2}$ $= \tfrac{1}{2}$}
\addplot[engblack, thick, dotted, domain=0:1] {0.7071};
\node[font=\scriptsize, anchor=south west, engblack] at (axis cs:0.02,0.71)
  {$V_{\mathrm{rms}}/V_{p} = 1/\sqrt{2}$};
\end{axis}
\end{tikzpicture}
\caption[Root-mean-square value of a sinusoid]{The root-mean-square
value of a sinusoid, the integral that converts a varying signal into
the single number that carries its power. The signal $v(t) = V_{p}
\sin(\omega t)$ (blue) swings between $\pm V_{p}$; its square (red) is
non-negative and oscillates at twice the frequency about its mean of
$\tfrac{1}{2}V_{p}^{2}$ (green dashed). The RMS value is the square
root of that mean, $V_{\mathrm{rms}} = V_{p}/\sqrt{2} \approx 0.707
V_{p}$ (black dotted), the DC voltage that would dissipate the same
power in a resistor. The factor $1/\sqrt{2}$ is the working content
of the integral $\frac{1}{T}\int_{0}^{T}\sin^{2}(\omega t)\dd t =
\tfrac{1}{2}$.}
\label{fig:vol02:ch06:rms}
\end{figure}
